\section{Exclusion of dense and nonoverlapping rare Bernoulli tails}\label{sec:rare-tails}

In this final section we show that the variational contact equations also
exclude infinite tails of rare two-point summands
\[
 X_j=h_j(p_j-B_j),\qquad B_j\sim\Bern(p_j),\qquad h_j>0,\qquad p_j\longrightarrow0.
\]

\subsection{Neumann-series probability identity at the maximizing threshold}

Set $\varphi=\phi(z)$ and define the right-hand contact defect
\begin{equation}\label{eq:right-defect}
 d(t)=F(z)+\varphi t-At^2+Ct^3-F(z+t)\ge0\qquad(t\ge0),
\end{equation}
which by monotonicity of $F$ satisfies
\begin{equation}\label{eq:right-defect-growth}
 d(t)\le\varphi t+|A|t^2+Ct^3,\qquad
 d(y)-d(x)\le K(y-x)\quad(0\le x\le y\le1).
\end{equation}
For a single summand $X_j=h(p-B)$ with $p<1/2$, $q=1-p$, and shift operator
$T_hf(t)=f(t+h)$, the distribution functions $\psi(t)=F(z+t)$ and
$g(t)=F_{S-X_j}(z-ph+t)$ satisfy $\psi=(qI+pT_h)g$. Inverting $qI+pT_h$ by
its absolutely convergent Neumann series gives, for any polynomially bounded
$f$ with $f(0)=0$,
\begin{equation}\label{eq:resolvent}
 \bigl[(T_h-I)(qI+pT_h)^{-1}f\bigr](0)
       =q^{-2}\sum_{k\ge1}(-p/q)^{k-1}f(kh).
\end{equation}
Since probability stationarity of $X_j$ gives
$[(T_h-I)(qI+pT_h)^{-1}\psi](0)=g(h)-g(0)=\varphi h-A(1-2p)h^2+C(1-2p)^3h^3$,
while $(T_h-I)(qI+pT_h)^{-1}$ maps $t,t^2,t^3$ at $0$ to
$h,(1-2p)h^2,(1-6pq)h^3$, subtracting yields the exact identity
\begin{equation}\label{eq:rare-score}
 \sum_{k\ge1}(-p/q)^{k-1}d(kh)
                   =-2Cp^2q^2(3-4p)h^3<0.
\end{equation}
In particular, moving the $k\ge2$ terms to the right-hand side and using
\eqref{eq:right-defect-growth} gives
\begin{equation}\label{eq:rare-defect-bound}
 d(h)\le\sum_{k\ge2}(p/q)^{k-1}d(kh)\le K_1ph
                      \qquad(p\le1/4,\ h\le1).
\end{equation}

\subsection{Exclusion of rare tails with dense gap spacings}

\begin{theorem}[Dense rare-gap exclusion]\label{thm:dense-rare-gaps}
Suppose a maximizing representation contains a sequence of distinct summands
$h_j(p_j-B_j)$ with $p_j\to0$ and $h_j\downarrow0$ satisfying, for some
$\nu>1$,
\begin{equation}\label{eq:actual-gap-mesh}
                 h_j-h_{j+1}=O(h_j^\nu).
\end{equation}
Then $F(z+t)=F(z)+\phi(z)t+O(t^{1+\alpha})$ as $t\downarrow0$ with
$\alpha=\min(1,\nu-1)$. Consequently, under the hypotheses of
Theorem~\ref{thm:weighted-variance}, no such sequence can exist if
\[
 \nu>1+\frac{\sqrt5-1}{2},
\]
or, when all but finitely many summands of the complete representation are
two-point, if
\[
 \nu>1+\frac{\sqrt{21}-3}{6}.
\]
\end{theorem}
\begin{proof}
For $1<\nu\le3$ and small $x>0$, let $s(x)$ be the largest $h_j\le x$, so
$0\le x-s(x)\le K_0x^\nu$ by \eqref{eq:actual-gap-mesh}. Combining
\eqref{eq:rare-defect-bound} at $h=s(x)$ with \eqref{eq:right-defect-growth}
shows that $H(x)=\sup_{0\le t\le x}d(t)$ satisfies
$H(x)\le c_0x^\nu+\sum_{k\ge2}r_0^{k-1}H(kx)$ for small $r_0>0$, which by a
barrier comparison with $K_2(x^\nu+x^3)$ implies $d(t)=O(t^\nu)$ as
$t\downarrow0$. Evaluating the phase-rounded array
\eqref{eq:weighted-phase-identity} at $z+b\ell$ ($b\asymp\sqrt{J\log J}$,
$\ell=J^{-a}$ with $\frac{1+\alpha}{2\alpha}<a<1+\theta$) then contradicts
maximality whenever $\alpha>1/(1+2\theta)$, which under the grouping rates
$\theta=(2s-1)/(6-s)$ ($s<\sqrt5-1$) and $\theta=(2s+3)/(6-s)$
($s<\sqrt{21}-3$) of Theorem~\ref{thm:weighted-variance} yields the stated
thresholds.
\end{proof}

\subsection{Exclusion of nonoverlapping rare tails}

\begin{theorem}[Nonoverlapping rare-tail exclusion]\label{thm:nonoverlapping-tail}
No maximizing representation can have the form
\[
 S=P+\sum_{j\ge1}h_j(p_j-B_j),
\]
where $P$ has finite support and is independent of the tail, $0<p_j\le1/2$
with $p_j\to0$, and for some constant $c>0$,
\[
 T_j:=\sum_{k>j}h_k\le h_j,
 \qquad h_{j+1}\ge c h_j>0
\]
for all sufficiently large $j$ (for example, $h_j=ar^j$ with $0<r\le1/2$ and
$p_j\to0$).
\end{theorem}
\begin{proof}
Put $M(t)=F(z+t)-F(z)$. By \eqref{eq:right-defect}--\eqref{eq:rare-defect-bound}
and $h_j\ge h_{j+1}+h_{j+2}\ge2h_{j+2}$, we have $M(t)\le K_3t$ and
\begin{equation}\label{eq:rare-annulus}
                  M(h_j)-M(h_{j+2})\ge\frac{\varphi h_j}{3}
\end{equation}
for all large $j$. For each atom $a\in\supp P$ of mass $w_a$ with
$y_a=a+\sum_jp_jh_j-z=\sum_jh_j\varepsilon_{a,j}\in\{\sum_jh_j\varepsilon_j\}$,
the prefix weight $v_{a,j}=w_a\prod_{k<j}p_k^{\varepsilon_{a,k}}(1-p_k)^{1-\varepsilon_{a,k}}$
at any index with $\varepsilon_{a,j}=1$ satisfies
\begin{equation}\label{eq:rare-prefix-cap}
 v_{a,j}(1-p_j)\le M(2h_j)\le2K_3h_j,
                         \qquad v_{a,j}\le4K_3h_j.
\end{equation}
Because $T_j\le h_j$ and $p_j\to0$, contributions to $M(h_i)-M(h_{i+2})$ from
first-difference digits $j<i$ are $o(h_i)$, while digits $>i+2$ have
displacement $\le T_{i+2}\le h_{i+2}$. Thus for each large $i$, some
$a\in\supp P$ and $j\in\{i,i+1,i+2\}$ must have $\varepsilon_{a,j}=1$ and
$v_{a,j}\ge[\varphi/(12|\supp P|)]h_j$, which forces $v_{a,k}/h_k\le4K_3c^{-(k-j)}p_j\to0$
for the next $L=3|\supp P|+3$ indices, contradicting the pigeonhole principle
over $|\supp P|+1$ disjoint three-index windows.
\end{proof}
